\subsection{Wall coatings as an effective boundary condition}
\label{subsec:wall_coatings}

%The previous sections considered the loss of spin-wave overlap due to atomic motion inside the optical interaction region. This description is sufficient when the optical beams occupy only a small fraction of the transverse cell area, because atoms can diffuse out of the mode before the cell boundary becomes the dominant limitation. In the compact geometry considered here, however, the situation changes. The cell diameter is only \(1~\mathrm{mm}\), while the control and signal beams have FWHM diameters of approximately \(0.8~\mathrm{mm}\) and \(0.3~\mathrm{mm}\), respectively. The optical mode therefore occupies a substantial fraction of the cell cross-section. Under these conditions, diffusion no longer occurs in an effectively unbounded transverse volume: atoms that carry part of the stored spin-wave coherence inevitably encounter the cell wall on the timescale relevant for the memory decay. A realistic interpretation of the simulated decay time \(\tau\) must therefore include not only transport through the buffer gas, but also the boundary condition imposed by the cell surface.

%During the residence time on the surface, the atom is exposed to interactions that are absent during free flight. Happer treats this type of relaxation as arising from fluctuating perturbations to the atomic Hamiltonian, which can include magnetic, electric, and spin-dependent interactions at the surface \cite{OP_Happer_1972}. For alkali atoms on uncoated glass, these perturbations are strong enough that wall collisions have historically been treated as highly depolarizing in optical-pumping experiments \cite{OP_Happer_1972}. The atom may therefore return mechanically to the vapor, but its ground-state coherence can be lost during the surface interaction.

The wall interaction must be separated into two conceptually different processes: mechanical reflection and internal-state survival. Mechanical reflection determines whether the atom remains inside the cell after reaching the boundary. Spin survival determines whether the atom still contributes coherently to the collective ground-state excitation after the collision. Happer describes relaxation in optically pumped vapors using the density matrix and treats many relaxation processes as the result of a weak fluctuating perturbation added to the atomic Hamiltonian,
     $
     \hat{H}(t) = \hat{H}_0 + \hat{V}(t),
    \label{eq:weak_fluctuation_hamiltonian}
    $
where \(\hat{H}_0\) is the unperturbed atomic Hamiltonian and \(\hat{V}(t)\) represents a random perturbation that varies from atom to atom and from collision to collision \cite{OP_Happer_1972}. In the case of wall relaxation, this perturbation is not applied continuously during free flight. Instead, it acts primarily when the atom interacts with the cell surface. Thus, an atom may elastically return to the vapor while its electronic or hyperfine coherence has already been randomized.

Upon striking the wall, an atom does not necessarily rebound immediately. It may stick to the surface for a finite dwell time before desorbing back into the vapor \cite{OP_Happer_1972,Aiacoavc_ChiEtAl_2020}. During this residence time, thermal agitation can cause the atom to hop between adsorption sites. Happer gives the average dwell time in the form
\begin{equation}
    \tau_{\mathrm{s}}
    =
    \tau_0
    \exp\!\left(\frac{E_{\mathrm{a}}}{k_{\mathrm{B}}T}\right),
    \label{eq:dwell_time}
\end{equation}
where \(\tau_0\) is a microscopic attempt time, \(E_{\mathrm{a}}\) is the adsorption energy, \(k_{\mathrm{B}}\) is Boltzmann's constant, and \(T\) is the wall temperature \cite{OP_Happer_1972}. While the atom is adsorbed, local magnetic fields, electric-field gradients, spin-orbit interactions, or nuclear spins in the substrate can act as random perturbations \cite{OP_Happer_1972}. These perturbations generate random phase shifts and spin relaxation. For an uncoated glass wall, this interaction is typically strong enough that a wall encounter can be treated, to a good approximation, as a fully depolarizing event for an alkali atom. In the absence of buffer gas, the relaxation time is then closely connected to the average flight time between wall collisions and therefore scales with the characteristic linear dimension of the cell \cite{OP_Happer_1972}. This scaling is precisely why wall effects become increasingly important as the vapor cell is miniaturized.

Anti-relaxation coatings reduce this surface-induced relaxation by separating the alkali atom from the strongly perturbing glass substrate. The coating provides a molecular layer in which the atom can adsorb and desorb with a smaller probability of spin randomization. If the perturbation accumulated during one surface encounter is weak, the spin state is only slightly changed after a single bounce and the atom can survive many wall bounces before losing its polarization. The coating can therefore be characterized by an average number of coherence-preserving collisions, denoted here by \(N_{\text{b}}\). Equivalently, one can describe the coating by a depolarization probability per bounce. This effective description is useful because the full microscopic surface dynamics are complex and material dependent.

Following the classification framework established by Chi et al.~\cite{Aiacoavc_ChiEtAl_2020}, the main coating families for alkali-vapor cells include straight-chain alkanes, alkenes, and organochlorosilanes. Straight-chain alkanes, such as paraffin, form ``wet'' films mainly through physical adsorption on the inner glass surface. These coatings were among the first successful anti-relaxation coatings and can allow alkali atoms to undergo many wall collisions before depolarization, with typical reported values around \(N_{\text{b}}\sim 10^4\). Their chemical structure is dominated by saturated C--C bonds. These bonds have relatively low polarizability and weak adsorption energies, which helps reduce spin-randomizing interactions during wall contact. However, paraffin-based coatings are limited by thermal stability and are generally unsuitable for high-temperature operation.

Alkene coatings are also physically adsorbed films, but they contain unsaturated C=C bonds. These bonds are more polarizable than saturated C--C bonds, and early intuition might suggest that this would be harmful for spin preservation. However, experiments discussed in the coating literature show that C=C bonds are not necessarily strongly depolarizing, and alkene coatings can reach very large bounce numbers, in some cases approaching \(N_{\text{b}}\sim 10^6\) \cite{Aiacoavc_ChiEtAl_2020}. Their main limitation is again temperature stability: their anti-relaxation performance can degrade at temperatures that are relevant for high-density alkali vapors. This makes them attractive from the point of view of spin survival, but less suitable when high optical depth requires elevated cell temperature.

Organochlorosilanes, such as octadecyltrichlorosilane (OTS), represent a different chemical strategy. Instead of forming only a physically adsorbed wet film, these coatings can form chemically bonded dry films. The chlorosilane end groups react with hydroxyl groups on the glass or silicon surface, producing siloxane bonds that anchor the organic layer to the substrate \cite{Aiacoavc_ChiEtAl_2020}. This chemical attachment improves thermal robustness and allows operation at substantially higher temperatures than typical paraffin or alkene coatings. The trade-off is that the spin survival is usually lower, with reported maximum values closer to \(N_{\text{b}}\sim 10^3\) than to the best alkene coatings. Thus, organochlorosilanes are not necessarily the best coatings in terms of bounce number, but they are important for compact warm-vapor memories because they remain usable in temperature regimes where high optical depth becomes accessible.

The coating literature also shows that the anti-relaxation mechanism is not completely settled. For paraffin, models based on adsorption, electron-randomization relaxation, spin-exchange relaxation, and relaxation associated with atoms in stems or reservoirs have been developed. Chi et al. emphasize that surface-science techniques are increasingly used to study coating composition and morphology, including methods such as atomic force microscopy, ellipsometry, X-ray diffraction, and contact-angle measurements \cite{Aiacoavc_ChiEtAl_2020}. These measurements are relevant because coating performance depends not only on the chemical family, but also on thickness, roughness, surface coverage, molecular ordering, adsorption energy, and the interaction between the coating and the alkali atom. For this reason, \(N_{\text{b}}\), atomic dwell time, hyperfine relaxation, and storage time are best understood as effective experimental characterizations of a complicated microscopic surface.

Temperature makes this trade-off central for warm-vapor quantum memories. Increasing the cell temperature raises the alkali vapor pressure and therefore the optical depth \(d\), which is beneficial for storage and retrieval efficiency. At the same time, a higher temperature increases the thermal velocity \(v_{\mathrm{th}}\), modifies the diffusion constant, increases the frequency with which atoms sample the cell boundary, and can degrade or modify the coating itself. Temperature therefore plays two opposite roles. It improves the optical interaction by increasing atomic density, but it can reduce the spin-wave lifetime through faster motion and stronger wall-related relaxation. This creates the central optimization problem for the compact cell: the memory should operate at a temperature high enough to provide sufficient optical depth, but not so high that wall-induced relaxation or coating degradation dominates the decay.

The present simulation does not attempt to resolve the microscopic chemistry of the coated surface. Instead, the coating is represented as an effective stochastic boundary condition. When an atom reaches the cylindrical wall or one of the end caps, its trajectory is reflected according to the chosen kinematic boundary rule. Independently of this mechanical reflection, its spin-wave contribution is multiplied by a survival factor. A simple parametrization is
\begin{equation}
    p_{\mathrm{survive}}
    =
    1 - \frac{1}{N_{\text{b}}},
    \label{eq:survival_probability}
\end{equation}
where \(N_{\text{b}}\) is the average number of wall collisions an atom can undergo before losing its coherence. This parameter condenses the surface material, adsorption dynamics, and microscopic relaxation mechanisms into a single quantity that can be varied numerically.

This abstraction also clarifies how coatings change the interpretation of the simulation results. In the low-survival limit, corresponding to bare glass or a poor coating, wall collisions act as an absorbing boundary for the spin wave. In that regime, increasing the buffer-gas pressure can improve the memory lifetime by slowing diffusion and reducing the probability that atoms reach the wall during the storage time. In the high-survival limit, corresponding to an effective anti-relaxation coating, atoms can collide with the wall and later re-enter the optical mode while still carrying coherence. This can produce long non-exponential tails in the retrieved coupling, because part of the ensemble temporarily leaves the high-overlap region and later returns.

The simulations therefore probe a regime that is not captured by the two simplest analytical limits. The limit \(N_{\text{b}}=0\) corresponds to immediate spin destruction at the wall, while \(N_{\text{b}}\rightarrow\infty\) corresponds to perfect coherence preservation. Real coated compact cells lie between these extremes. In particular, when the cell radius, beam waist, and diffusion length are comparable, the retrieval decay depends on the combined effect of beam geometry, buffer gas, temperature, and wall survival. The numerical results show that increasing \(N_{\text{b}}\) strongly improves the decay for poor coatings, but that the curves tend to converge once \(N_{\text{b}}\sim 10^3\text{ to }10^4\). Above this range, wall-induced depolarization is no longer the dominant loss channel. The remaining reduction of coupling is then mainly set by diffusion out of the signal and control mode volume, finite cell geometry, and the spatial overlap of the spin wave with the readout mode. This is why the coating model is essential for interpreting the simulated \(\tau\)-versus-temperature trends in compact vapor cells.
